\documentclass{article}
\usepackage{graphicx} % Required for inserting images
\usepackage{minted}
\usepackage[a4paper, margin=2cm]{geometry}
\usepackage{tcolorbox}
\usepackage{hyperref}
\usepackage{xcolor}

% This sets the global style for all minted environments
\setminted{
    bgcolor=gray!10, % A very subtle 10% gray
    linenos=true,    % Optional: adds line numbers
}

\title{Tutorato Programmazione 2}
\author{Michele Porcellato, Giacomo Vettore}
\date{6 Maggio 2026}

\begin{document}

\maketitle

\section{Soluzioni Teoria}

Dovete valutare i seguenti codici ed indicare l’output oppure l’errore. In caso di errore, indicate a che
riga è l’errore e, a grandi linee, di che tipo di errore si tratta.

\begin{minted}{java}
import java.util.*;
public class Test {
    public static void main(String[] args) {
        String[] a = {"limone", "ananas", "mango", "lime"};
        Set<Integer> s = new HashSet<>();
        List<Integer> l = new ArrayList<>();
        for(String i: a) {
            s.add(i.length());
            l.add(i.length());
        }
        System.out.println(s.size() + l.size());
    }
}
\end{minted}

\begin{tcolorbox}[
    width=\textwidth,
    height=5cm,
]
7
\end{tcolorbox}

\begin{minted}{java}
public class Test {
    public static void main(String[] args) {
        I i = new C(4);
        System.out.println(i.m(2));
    }
}
interface I {
    int m(int z);
}
class A implements I {
    int x;
    A(int x) { this.x = x + 1; }
    public int m(int z) { return x + z; }
}
class B extends A { 
    B(int x) { super(++x); }
    public int m(int z) { return x * z; }
}
class C extends B {
    C(int x) { super(++x); }
}
\end{minted}

\begin{tcolorbox}[
    width=\textwidth,
    height=5cm,
]
14
\end{tcolorbox}

\begin{minted}{java}
import java.util.ArrayList;
import java.util.List;

public class Test {
    public static void M(List<String> lista) {
        System.out.println("Elaborazione lista di stringhe...");
        for (String s : lista) {
            System.out.println("- " + s);
        }
    }
    
    public static void M(List<Integer> lista) {
        System.out.println("Elaborazione lista di interi...");
        for (Integer i : lista) {
            System.out.println("- " + i);
        }
    }

    public static void main(String[] args) {
        List<String> A = new ArrayList<>();
        A.add("Alice");
        A.add("Bob");

        List<Integer> B = new ArrayList<>();
        B.add(10);
        B.add(20);

        M(A);
    }
}
\end{minted}

\begin{tcolorbox}[
    width=\textwidth,
    height=5cm
]
\textbf{Errore:} Compile-time \\
error: name clash: M(List) and M(List) have the same erasure \\
\# \textit{Commento:} Visto che i Generics di Java sono implementati con type-erasure, il compilatore, dopo aver fatto type-checking, non è in grado di distinguere i due metodi (hanno di fatto la stessa firma). \\
Riga 12
\end{tcolorbox}

\end{document}